\documentclass[11pt]{article}
\usepackage[margin=1in]{geometry}
\usepackage{amsmath,amssymb,booktabs,hyperref,graphicx}
\title{Surrogate Replication of the Correlation-Driven $q=1/2$ Charge Order\\
in the Kagome Superconductor CeRu$_3$Si$_2$\\
\large (Gerguri \textit{et al.}, arXiv:2603.27408, 2026)}
\author{Automated replication (Hermes subagent) --- tight-binding + mean-field surrogate}
\date{\today}
\begin{document}
\maketitle

\begin{abstract}
We test the headline claim of Gerguri \textit{et al.} (2026): that DFT+U reproduces
the experimentally \emph{dominant} $q=1/2$ (Pmma) charge order in CeRu$_3$Si$_2$
only for a Ce-$4f$ Hubbard interaction $U>6$~eV, with the weaker $q=1/3$ (Imma)
order remaining nearly degenerate, while treating Ce-$4f$ as core fails to
stabilize $q=1/2$. Full DFT+U is scoped out; instead we build a from-scratch
kagome tight-binding + Hartree mean-field \emph{surrogate} (Ru kagome layer
hybridized with a heavy Ce-$4f$ flat level) and compare the charge-order
susceptibility $\chi_q$ at $q=1/2,\,1/3,\,1/4$ versus $U$. The surrogate
reproduces the paper's mechanism qualitatively: $q=1/3$ wins at small $U$,
a crossover to $q=1/2$ occurs at $U\!\approx\!5$ (paper: $6$~eV) with the two
nearly degenerate near the crossover, and the Ce-$4f$-as-core control yields a
$q=1/4$ (CO$^*$) winner with $q=1/3$ suppressed --- matching the paper's reported
failure mode. Verdict: \textbf{PARTIAL} (qualitative-only; no first-principles
total energies).
\end{abstract}

\section{Paper claim and target}
The parent phase is P6/mmm kagome. Synchrotron XRD reveals a dominant $q=1/2$
charge order (describable in Pmma) plus a weaker $q=1/3$ component (Imma, CO-II),
persisting to 300~K. First-principles DFT shows coexisting Ru $d$-orbital kagome
flat bands and heavy-fermion Ce$^{4+}$ $4f$ flat bands (4f$^0$, above $E_F$), with
multiple imaginary phonon modes in the $q_z=1/2$ plane. Crucially, with Ce-$4f$ as
valence: at $U=0$ the $q=1/3$ order is favored (ground state actually a non-CO
distortion); the two orders are \emph{nearly degenerate at $U=6$~eV}; and for
$U>6$~eV the $q=1/2$ order becomes the lowest-energy ground state. With Ce-$4f$ as
core, $q=1/3$ is suppressed and $q=1/2$ is not stabilized (a $q=1/4$ CO$^*$
appears).

\section{Method (surrogate)}
We diagonalize a real-space kagome supercell: 3 Ru sites/cell (nearest-neighbor
hopping $t=1$) plus one Ce-$4f$ level/cell (onsite $\varepsilon_f$, Ru--$f$
hybridization $t_f=0.30$). The Hubbard $U$ is mapped to the $4f$ flat-band
position, $\varepsilon_f=\varepsilon_{f,0}+k_U U$ ($\varepsilon_{f,0}=0.15$,
$k_U=0.10$), encoding that stronger correlations push the $4f^0$ band further above
$E_F$. A commensurate charge order at $q=1/n$ is imposed as an onsite Ru modulation
$\delta\cos(2\pi x/n)$. The competing order is selected by the mean-field CDW
susceptibility (Landau curvature)
\begin{equation}
\chi_q = -\frac{E(+\delta)+E(-\delta)-2E(0)}{\delta^2},\qquad
E=\tfrac{1}{N}\!\!\sum_{\text{occ}}\!\varepsilon_i ,
\end{equation}
evaluated at filling $0.62$ (near the Ru kagome M-point van Hove condition). The
largest $\chi_q$ is the favored order. Geometry and tight-binding conventions are
adapted from \texttt{loop\_current\_kagome\_kernel.py}; the occupied-energy /
density-matrix machinery follows \texttt{loop\_current\_meanfield\_kernel.py}
(both credited).

\section{Results}
\begin{table}[h]\centering
\begin{tabular}{rccccl}
\toprule
$U$ (eV) & $\varepsilon_f/t$ & $\chi_{1/2}$ & $\chi_{1/3}$ & $\chi_{1/4}$ & ground state\\
\midrule
0 & 0.15 & 0.389 & \textbf{0.448} & 0.292 & $q=1/3$\\
2 & 0.35 & 0.350 & \textbf{0.401} & 0.225 & $q=1/3$\\
4 & 0.55 & 0.313 & \textbf{0.324} & 0.179 & $q=1/3$\\
5 & 0.65 & \textbf{0.297} & 0.278 & 0.162 & $q=1/2$\\
6 & 0.75 & \textbf{0.284} & 0.233 & 0.149 & $q=1/2$\\
7 & 0.85 & \textbf{0.274} & 0.194 & 0.140 & $q=1/2$\\
8 & 0.95 & \textbf{0.268} & 0.165 & 0.133 & $q=1/2$\\
9 & 1.05 & \textbf{0.265} & 0.145 & 0.129 & $q=1/2$\\
\midrule
\multicolumn{2}{l}{Ce-$4f$ as core} & 0.285 & 0.109 & \textbf{0.472} & $q=1/4$ (CO$^*$)\\
\bottomrule
\end{tabular}
\caption{Charge-order susceptibility $\chi_q$ vs Hubbard $U$. $q=1/3$ wins at
small $U$; crossover to $q=1/2$ at $U\!\approx\!5$; near-degeneracy at
$U\!\approx\!4$--$5$. The $f$-as-core control gives a $q=1/4$ CO$^*$ winner with
$q=1/3$ strongly suppressed --- reproducing the paper's failure mode.}
\end{table}

\section{Comparison and verdict}
\begin{itemize}
\item \textbf{Reproduced (qualitative):} (i) $q=1/3$ favored at $U=0$; (ii) a
correlation-driven crossover to a $q=1/2$ ground state at large $U$; (iii)
near-degeneracy of $q=1/2$ and $q=1/3$ around the crossover; (iv) the $f$-as-core
control fails to stabilize $q=1/2$ and instead favors $q=1/4$ (CO$^*$) with
$q=1/3$ suppressed --- three independent qualitative agreements with the paper.
\item \textbf{Not reproduced (quantitative):} the surrogate crossover is at
$U\!\approx\!5$ (arbitrary $t$-units mapped to eV), not a first-principles $6$~eV;
no true total energies; phonon channel omitted; no self-consistency.
\item \textbf{Verdict:} PARTIAL --- gap: no DFT+U total-energy comparison
(scoped out); surrogate is qualitative only.
\end{itemize}
Coverage: 7/10. Agreement: 7/10.

\section{Kernel credit}
Physics built on the shared TEXTURES-100 kernels
\texttt{loop\_current\_kagome\_kernel.py} (kagome geometry, primitive vectors,
NN Bloch/real-space construction) and
\texttt{loop\_current\_meanfield\_kernel.py} (occupied-density / band-energy
evaluation). See \texttt{report/evidence/}.
\end{document}
